\section{Restoring generation and two transverse directions}
\label{sec:generation}

The preceding construction gives a substitution $Z(s_1,s_2)$ inside
the centralizer of $A$ in a larger central division algebra.  It
preserves the old polynomial identities and algebraizes the chosen
series.  To complete Theorem~\ref{thm:lifting}, we must
make the new constant tuple generate the larger algebra and retain
two independent directions modulo inner derivations. We first use the
original transverse directions to prove generation of the generic
diagonal product and to control the arrows between its blocks.
Translating the first free letter by a new central parameter then joins
the blocks into a full matrix algebra. A single rank calculation shows
that both directions remain independent modulo inner derivations.

Throughout this section, brackets with a tuple are taken componentwise.
Thus, for $C=(C_1,\ldots,C_u)$ in an algebra $B$, write
\[
 J(C)=\{([T,C_1],\ldots,[T,C_u]):T\in B\}.
\]
Polynomial relations on $X_0$ detect the two input directions to first
order.

\begin{lemma}[Derivatives of relations]
\label{lem:derivatives-of-relations}
Let $A$ be a finite-dimensional central simple algebra over a field
$k$ of characteristic zero.  Let
$X\in A\langle\!\langle z_1,z_2\rangle\!\rangle^u$ have constant
tuple $X_0$ generating $A$, and suppose its two linear coefficients
$X_1,X_2$ are independent modulo $J(X_0)$.  Define
\[
 \theta:k\langle x_1,\ldots,x_u\rangle\longrightarrow A,
 \qquad \theta(f)=f(X_0),\qquad I=\ker\theta.
\]
Then the map taking $f\in I$ to the coefficients of $z_1,z_2$ in
$f(X)$ is onto $A^2$.
\end{lemma}

\begin{proof}
Denote these two coefficient maps by $\delta_1,\delta_2$.  Each
satisfies
\[
 \delta_j(fg)=\theta(f)\delta_j(g)+\delta_j(f)\theta(g).
\]
Since $\theta$ is onto and $I$ is a two-sided ideal, the image
$M=(\delta_1,\delta_2)(I)$ is an $A$-bimodule in $A^2$.

Left and right multiplication generate $\operatorname{End}_k(A)$;
see \cite[Chapter~IV, Corollary~2.9]{MilneCFT}.
Indeed, after splitting $A$, multiplication by matrix units on the
left and right gives every matrix unit on the vector space $A$;
the resulting isomorphism
$A\otimes_k A^{\mathrm{op}}\simeq\operatorname{End}_k(A)$ descends
by finite-dimensional linear algebra.  Applying these endomorphism
matrix units to $M\subseteq A\otimes_k k^2$ shows that
$M=A\otimes_k V$ for a subspace $V\subseteq k^2$.
If $V\ne k^2$, some nonzero $(a,b)\in k^2$ annihilates $V$.
The map $a\delta_1+b\delta_2$ then vanishes on $I$ and descends to
a derivation of $A$.

Every $k$-linear derivation of $A$ is inner
\cite[Theorem~8]{Jacobson1937Derivations}. For completeness, after
splitting $A=M_d(E)$, a derivation $\partial$ is implemented by
\[
 T=\sum_i\partial(E_{i1})E_{1i}.
\]
Differentiating $E_{kl}E_{i1}=\delta_{li}E_{k1}$, multiplying by
$E_{1i}$, and summing gives $\partial(E_{kl})=[T,E_{kl}]$.
The equations for an implementing element on a fixed $k$-basis of
$A$ form a linear system over $k$.  Solvability after scalar
extension implies solvability over $k$.
It follows that $aX_1+bX_2\in J(X_0)$, a contradiction.
\end{proof}

In particular, there are $f_1,f_2\in I$ whose derivative pairs are
$(1,0)$ and $(0,1)$.  These relations recover the two generic input
matrices from small evaluations of the tuple.

\begin{lemma}[Generic blocks]
\label{lem:generic-generating-blocks}
Under the hypotheses of Lemma~\ref{lem:derivatives-of-relations},
write $d=\deg A$ and split $A$ into $d\times d$ matrices.  Fix a
positive integer $b$ and finitely many independent small generic input pairs
\[
 P_i=\varepsilon P_i^{\circ},\qquad 1\le i\le r,
\]
of $b\times b$ matrices, with jointly independent entries in the
$P_i^{\circ}$.  Work over a common splitting Laurent field.
Then each tuple $C_i=X(P_i)$ generates $M_{db}$, and the diagonal
tuple $C=\operatorname{diag}(C_i)$ generates the full product of
these matrix algebras.

For distinct $i,j$, the output derivative of a generic input arrow
$U=(U_1,U_2)$ from block $j$ to block $i$ is nonzero modulo
\[
 \{C_iT-TC_j:T\in M_{db}\}.
\]
More precisely, the image of the arrow derivatives in this quotient
has dimension at least $b^2$.
\end{lemma}

\begin{proof}
Let $E$ contain the splitting data and independent input coordinates.
Consider integral elements of the $E((\varepsilon))$-algebra
generated by $X(P_i)$.  Their reductions modulo $\varepsilon$
contain $M_d(E)\otimes1$, because $X_0$ generates the old algebra.
Lemma~\ref{lem:derivatives-of-relations} also gives
\[
 \varepsilon^{-1}f_j(X(P_i))
 \equiv I_d\otimes P_{i,j}^{\circ}\pmod\varepsilon.
\]
A generic pair generates $M_b(E)$.  To see the relevant determinant
is nonzero, specialize one matrix to a diagonal matrix with distinct
entries and the other to a matrix with every entry nonzero.
Interpolation and sandwiching give all matrix units.  Thus the
reductions generate $M_d(E)\otimes M_b(E)=M_{db}(E)$.
Choose finitely many integral algebra elements whose reductions
form a matrix basis.  Their determinant is a unit, so they remain
independent over $E((\varepsilon))$, proving full generation.

The representations defined by different $C_i$ have distinct kernels.
Indeed,
\[
 \operatorname{tr}f_1(C_i)
  =\varepsilon d\operatorname{tr}P_{i,1}^{\circ}
    +O(\varepsilon^2),
\]
and the leading coefficients differ.  Equal kernels would induce an
isomorphism of the full matrix images carrying every polynomial
evaluation in one to the other.  Such an isomorphism preserves trace:
the trace of left multiplication is matrix size times matrix trace.
This contradicts the display.
The kernels are therefore distinct maximal two-sided ideals.
Evaluation is onto their finite product.  Explicitly, for each pair
$i\ne j$ choose a polynomial vanishing in block $j$ and equal to
the identity in block $i$.  The product over $j\ne i$ vanishes in
all other blocks and is the identity in block $i$; multiplying by
a lift of any matrix gives that product coordinate.

For the last assertion, compare the inner columns and arrow columns
at $P_i=P_j=0$.  The inner image is
$J(X_0)\otimes M_b$, of dimension $(d^2-1)b^2$.  The arrow derivative
is
\[
 X_1\otimes U_1+X_2\otimes U_2.
\]
The original transverse independence implies that the combined map
has rank $(d^2+1)b^2$ at this reduction.  Its generic rank is at
least this number.  At a generic pair, the inner map has rank at
most $(db)^2=d^2b^2$; in fact it is injective, since a nonzero
intertwiner between the two simple representations would identify
their kernels.  Subtracting this generic inner rank leaves at least
$b^2$ independent arrow directions in the quotient.  A generic
fixed arrow consequently has nonzero quotient class.
\end{proof}

All determinant conditions in Lemma~\ref{lem:generic-generating-blocks}
hold on the same independent block and arrow coordinates used in
the division realization.  Only finitely many blocks and arrows are
involved, so their nonzero minors can be imposed simultaneously.
In particular, the argument does not require replacing that realization
by a different choice of data.

We next isolate the translation argument.  Its hypotheses concern
only the constant and linear coefficients; higher-order corrections
may be arbitrary.

\begin{lemma}[A translation preserving two directions]
\label{lem:translation-two-directions}
Let $E$ be a field of characteristic zero, let $n\ge1$ be an integer,
let $m=11$, and index
the blocks of $M_{mn}(E)$ by $\mathbb Z/m\mathbb Z$.
Suppose
\[
 Y(s_1,s_2)\in M_{mn}(E)
       \langle\!\langle s_1,s_2\rangle\!\rangle^u
\]
has constant tuple $C=\operatorname{diag}(C_i)$ generating
$\prod_iM_n(E)$.  Suppose its two linear coefficients have the
following supports:
\[
 (Y_1)_{ij}=0\ \text{unless }i-j\in\{1,2\},
 \qquad
 (Y_2)_{ij}=0\ \text{unless }i-j\in\{3,4\}.
\]
Assume every displayed edge coefficient is nonzero modulo the
cross-inner image $\{C_iT-TC_j:T\in M_n(E)\}$.
For a new central Laurent parameter $\gamma$, define
\[
 Y'(s'_1,s'_2)=Y(\gamma+s'_1,s'_2).
\]
Then $Y'_0=Y(\gamma,0)$ generates $M_{mn}(E((\gamma)))$, and
$Y'_1,Y'_2$ are independent modulo $J(Y'_0)$.
\end{lemma}

\begin{proof}
The substitution is well-defined.  In the coefficient of a fixed
new word, a contribution from an old word of length $r$ contains
$r$ minus the new word length factors of $\gamma$.  Thus only
finitely many terms contribute at each $\gamma$ degree.  Coefficients
converge in $E[[\gamma]]$, and the substitution respects products.

Let $B=\prod_iM_n(E)$ and let $I_C$ be the kernel of evaluation of
the free algebra at $C$.  First-letter differentiation on $I_C$
has image a $B$-bimodule in $M_{mn}(E)$.  Its projection onto every
first-color edge is nonzero.  Otherwise the derivative into that
cross bimodule would descend to a derivation $\delta:B\to M_n(E)$,
where the left and right actions are through the distinct endpoint
factors $i,j$.  Such a derivation is inner.  In fact, for the central
idempotent $e_i$ of factor $i$, differentiating $e_i a=a e_i$ gives
\[
 \delta(a)=a_iT-Ta_j,\qquad T=\delta(e_i).
\]
Applied to the generators, this contradicts the transverse edge
hypothesis.  Since each full rectangular block is a simple bimodule
for its two matrix factors, the derivative image contains the full
block on every first-color edge.

Divide the corresponding evaluations of elements of $I_C$ at
$Y(\gamma,0)$ by $\gamma$.  Their reductions supply all these edge
blocks.  Ordinary evaluations already supply the full diagonal
product.  The offset-one edges form a directed cycle through all
vertices, so products supply every matrix block.  A finite set of
such elements with independent reductions therefore proves that
$Y(\gamma,0)$ generates the full matrix algebra.

It remains to check the two directions.  Put $N=mn$ and
$J_0=J(C)$.  Since $C$ generates $B$, its commutant consists of the
$m$ block scalars.  Hence $\dim J_0=N^2-m$.  Choose $N^2-m$ inner
columns with independent reductions spanning $J_0$.  For a block
scalar $c=\operatorname{diag}(c_i I_n)$, the additional column
\[
 \gamma^{-1}[c,Y(\gamma,0)]
\]
is integral and reduces to $[c,Y_1]$.  These columns have rank
$m-1$ modulo $J_0$.  Indeed, $[c,Y_1]\in J_0$ forces $c_i=c_j$
on each offset-one edge, by its transverse hypothesis, and the
cycle forces all $c_i$ to agree.

The two shifted linear coefficients reduce to $Y_1,Y_2$.  Suppose
\[
 aY_1+bY_2\in J_0+\{[c,Y_1]:c\text{ is block scalar}\}.
\]
At a second-color edge, only $bY_2$ contributes outside $J_0$,
so $b=0$.  On the directed offset-one cycle, the relation then
forces $a=c_i-c_j$ on each edge.  Summing gives $ma=0$, whence
$a=0$ in characteristic zero.  Thus these two reductions are
independent modulo the enlarged inner-column space.

We have found $N^2-m+(m-1)+2=N^2+1$ independent reductions among
inner columns and the two tangents.  They lift to independent
columns over $E((\gamma))$.  Full generation of $Y'_0$ implies
$\dim J(Y'_0)=N^2-1$, so the two tangent classes are independent.
\end{proof}

Apply the lemma to $Y=X(Z)$.  Lemma~\ref{lem:generic-generating-blocks}
supplies its diagonal-product and edge hypotheses, while the
eleven-block construction supplies precisely the two color supports.
Consequently the translated tuple generates after splitting and has
two independent tangent classes.  These are finite rank conditions,
so they also hold over the division algebra's own center.
The central Laurent extension by $\gamma$ preserves division.
Both the old polynomial identities and the new algebraic relation
are preserved by the convergent substitution.

\subsection{Returning to a countable center}

The complete Laurent fields used so far need not be countable.
We finish by retaining only the scalar data required by the result.
Use the notation of Section~\ref{sec:division-realization}:
$F_0=H_0(a_*)$ is the construction field, $D_0$ is the central
division algebra over it, and
\[
 L=F_0((\varepsilon))((\gamma))
\]
is the full shifted Laurent center. The field $F_0$ is countable:
$L_0$ is generated by finitely many indeterminates over the countable
field $k(\alpha,\beta)$, and $H_0$ is a subfield of $L_0$.
The algebra $D_0$ has degree $dbm$, and $D_0\otimes_{F_0}L$
is the central division algebra just constructed.

Choose an $F_0$-basis of $D_0$. Let $k'\subseteq L$ be the field
generated by $F_0$, the coefficients of the monic polynomial $q$
(taken to be $q_0$ when $q_0(w)=0$),
and every scalar coordinate in this basis of every coefficient
of the translated series $X'$ and $w'$. There are countably many
free words and finitely many coordinates per coefficient, so $k'$
is countable.  Put
\[
 A'=D_0\otimes_{F_0} k'.
\]
This is central simple of degree $dbm$.  Its scalar extension to
$L$ is division, so $A'$ embeds in a division ring and has no zero
divisors.  A finite-dimensional algebra without zero divisors is
a division algebra.  Thus $A'/k'$ has the required central division
property and strictly larger degree.

All coefficients of $X',w'$ and $q$ lie in this smaller model.
The identities descend by its injective scalar extension to $L$.
Generation and transverse independence descend as well, since
the ranks of the finite coordinate matrices defining them do not
change under field extension.  This proves every conclusion of
Theorem~\ref{thm:lifting}.

In the unconstrained case $q_0(w)=0$, the same argument applies
without the spectral corrections: use the same eleven-block
generic constant and colored linear terms, retain the relation
$q_0$, and carry out the translation and countable restriction
above.  Thus that case also increases degree while preserving
generation and both directions.
